% Nouveau EGA V (cote 25): archive cote 25, site page 179 (PDF page 180).
% typescript. Draft compared: none.
% AI-transcribed by Claude Opus 5.5 (claude-opus-5-5, adaptive thinking, effort high), Message Batches API, batch ega25b (msgbatch_01A7n5YuXubUPNK1zsepi8RG), 2026-11-13; not yet checked by eye.
\scanpage{25/179}
\mspage{\hand{6}}

\[ \underline{Q}(\underline{E}) =W(\underline{Sym}^2(\underline{E})) \]
et définissons de même $\underline{Q}(\underline{E}_o)$; par l'opération de ``transport de
structure'' on aura donc un morphisme $u \hand{\rightsquigarrow} \underline{Q}(u)$ :
\[ \underline{Isom}(\underline{E}_o,\underline{E}) \hand{\longrightarrow} \underline{\unsure{I|O}som}(\underline{Q}(\underline{E}_o),\underline{Q}(\underline{E})), \]
et, utilisant la section $\unsure{q|Q}_o$ de $\underline{Q}(\underline{E}_o)$ correspondant à $Q_o$ , on trouve un
morphisme $u \hand{\rightsquigarrow} \underline{Q}(u)(q_o)$:
\[ (*) \qquad \underline{Isom}(\underline{E}_o,\underline{E}) \hand{\longrightarrow} \underline{Q}(\underline{E}) \quad . \]
D'autre part, la forme Q \xout{sur E} correspond à une section q de $\underline{Q}(\underline{E})$ sur S,
et P n'est autre que l'image inverse de q(S) par (*), comme il résulte
\margin{\unsure{\illegible{?}}}
% NOTE: Faint bluish oblique marginal traces; no secure word or word count, and their status as front writing or show-through is unresolved.
trivialement des définitions, (ce qui, d'ailleurs, établit en passant la
représentabilité annoncée de P \unsure{c|x}omme préschéma affine et de \xout{type}\inserted[typed]{prés} finie
sur S). Introduiso\unsure{n|j}s maintenant l'ouvert \xout{$\underline{Q}$(x}
\[ \underline{Q}(\underline{E})^* = \underline{Q}(\underline{E})_{\overset{\hand{\sim}}{d}} \]
de $\underline{Q}(\underline{E})$ , correspondant aux formes quadratiques de discriminant corrigé
inversible (représentan\xout{x}t le foncteur $S' \hand{\rightsquigarrow}$ ensemble des \xout{formes}
sections de $\underline{Sym}^2(\underline{E}_{S'})$ à discriminant corrigé inversible). En fait,
par transport de structure on a $\underline{Isom}(\underline{E}_o,\underline{E}) \hand{\rightarrow} \underline{Isom}(\underline{Q}(\underline{E}_o)^*,\underline{Q}(\underline{E})^*)$, et
comme $q_o$ est une section de $\underline{Q}(\underline{E}_o)^*$ (puisque le discriminant \unsure{m|c}o\unsure{d|r}rigé de
la forme type $Q_o$ vaut 1), le morphisme (*) se factorise en fait en un
morphisme
\[ (*\,*) \qquad \underline{Isom}(\underline{E}_o,\underline{E}) \hand{\rightarrow} \underline{Q}(\underline{E})^* , \]
qui est évidemment de présentation finie, puisque les deux membres sont de
présentation finie sur S. Il suffit maintenant de prouver la
\xout{Lemme} \uline{Proposition} 16.11. Avec les notations précédentes (à rappeler),
le morphisme (* *) est \uline{fidèlement plat} .

Il s'ensui\deleted{t}\inserted[typed]{vra} que si Q satisfait à (ii), i.e. \inserted{si} q est en fait une sectio
% NOTE: The word is clipped after sectio at the right edge; no final n is visible.
de $\underline{Q}(\underline{E})^*$\tikz[baseline=-0.5ex,x=1em,y=1em]{\draw[handwritten,->] (0,0) .. controls (.6,0) and (.6,.5) .. (.6,.9);} , alors P (déduit de (**) par \inserted{le} changement de base q) est encore

\mspage{\hand{179}}
